\section{Probabilités conditionnelles}
\label{sec:conditionnement}

Prédire $Y$ à partir de $X$, c'est raisonner « sachant $X$ ». Cette section donne les deux outils du chapitre~1 : la formule de Bayes et la moyenne conditionnelle.

\subsection{Probabilité conditionnelle et formule de Bayes}

\begin{definition}[Probabilité conditionnelle]
Pour deux événements $A$ et $B$ avec $\P(B)>0$, la probabilité de $A$ \emph{sachant} $B$ est $\P(A\mid B)=\dfrac{\P(A\text{ et }B)}{\P(B)}$.
\end{definition}

\begin{proposition}[Probabilités totales et formule de Bayes]
\label{prop:bayes}
Si $B_1,\dots,B_K$ forment une partition (un et un seul se produit), de probabilités non nulles :
\[
  \P(A)=\sum_{k=1}^K\P(A\mid B_k)\,\P(B_k)
  \qquad\text{et}\qquad
  \P(B_j\mid A)=\frac{\P(A\mid B_j)\,\P(B_j)}{\sum_k\P(A\mid B_k)\,\P(B_k)} .
\]
\end{proposition}
\begin{preuve}
$A$ est la réunion des événements disjoints « $A$ et $B_k$ », de probabilités $\P(A\mid B_k)\P(B_k)$ : on additionne.
Puis $\P(B_j\mid A)=\P(A\text{ et }B_j)/\P(A)$, et l'on remplace $\P(A)$ par la première formule.
\end{preuve}

\begin{exemple}[Test de dépistage]
Une maladie touche $1\,\%$ d'une population. Un test est positif chez $95\,\%$ des malades et négatif chez $90\,\%$ des non-malades. Pour une personne testée positive :
\[
  \P(M\mid+)=\frac{0{,}95\times0{,}01}{0{,}95\times0{,}01+0{,}10\times0{,}99}=\frac{0{,}0095}{0{,}1085}\approx0{,}088 .
\]
Moins de $9\,\%$ des positifs sont malades : la rareté de la maladie pèse lourd. C'est le raisonnement du classifieur de Bayes (chapitre~1).
\end{exemple}

\subsection{Moyenne conditionnelle}

\begin{definition}[Loi et espérance conditionnelles, cas discret]
Si $X$ et $Y$ sont discrètes et $\P(X=x)>0$, la \emph{loi de $Y$ sachant $X=x$} est donnée par les probabilités $\P(Y=y\mid X=x)$, et
\[
  \E[Y\mid X=x]=\sum_yy\;\P(Y=y\mid X=x)
\]
est la \emph{moyenne de $Y$ sachant $X=x$} : la moyenne de $Y$ dans le groupe des individus pour lesquels $X=x$.
\end{definition}

Quand $Y$ ne prend que les valeurs $0$ et $1$, $\E[Y\mid X=x]=\P(Y=1\mid X=x)$ : une moyenne de $0$ et de $1$ est une proportion.
Le cas continu se définit de la même façon avec des densités ; nous l'admettrons.

\begin{proposition}[Formule de l'espérance totale]
\label{prop:tour}
$\E[Y]=\sum_x\E[Y\mid X=x]\;\P(X=x)$ : la moyenne générale est la moyenne pondérée des moyennes de groupes.
\end{proposition}
\begin{preuve}
$\sum_x\E[Y\mid X=x]\P(X=x)=\sum_x\sum_yy\,\P(Y=y\mid X=x)\P(X=x)=\sum_y y\sum_x\P(Y=y,X=x)=\sum_yy\,\P(Y=y)=\E[Y]$.
\end{preuve}

\begin{exemple}
Parmi les femmes d'une région, $40\,\%$ vivent en ville, où la moyenne est de $2{,}5$ enfants, et $60\,\%$ en milieu rural, avec $4{,}5$ enfants en moyenne.
La moyenne régionale est $0{,}4\times2{,}5+0{,}6\times4{,}5=3{,}7$ enfants.
Pour prédire le nombre d'enfants d'une femme dont on connaît le milieu, on utilisera $2{,}5$ ou $4{,}5$ plutôt que $3{,}7$ (chapitre~1).
\end{exemple}

\subsection*{Exercices}

\begin{exercice}[Prévalence et dépistage]
\label{ex:prevalence}
Reprendre le test de dépistage avec une prévalence de $10\,\%$. Que devient $\P(M\mid+)$ ? Commenter.
\end{exercice}

\begin{exercice}[Moyennes conditionnelles]
\label{ex:tableau}
Le couple $(X,Y)$ a la loi suivante.
\begin{center}
\begin{tabular}{c|ccc}
 & $Y=0$ & $Y=1$ & $Y=2$ \\ \hline
$X=0$ & $0{,}2$ & $0{,}1$ & $0{,}1$ \\
$X=1$ & $0{,}1$ & $0{,}2$ & $0{,}3$
\end{tabular}
\end{center}
Calculer $\E[Y\mid X=0]$ et $\E[Y\mid X=1]$, puis vérifier la formule de l'espérance totale.
\end{exercice}
